\documentclass[10pt,letterpaper]{article}

\usepackage[margin=1in]{geometry}
\usepackage{times}
\usepackage{amsmath,amssymb}
\usepackage{graphicx}
\usepackage{booktabs}
\usepackage{multirow}
\usepackage{xcolor}
\usepackage{hyperref}
\usepackage{caption}
\usepackage{subcaption}
\usepackage{natbib}
\usepackage{enumitem}

\hypersetup{colorlinks=true,linkcolor=blue,citecolor=blue,urlcolor=blue}

\title{Vision Encoders, Not Language Models, Drive Action-Relevant\\Representations in Vision-Language-Action Models}

\author{
  Anonymous Authors\\
  Conference on Robot Learning (CoRL) 2026 Submission
}

\date{}

\begin{document}
\maketitle

% ============================================================
\begin{abstract}
% ============================================================
Vision-Language-Action models (VLAs) fine-tune large language models on robot demonstration data to produce actions, yet the origin of action-relevant representational structure within these models remains unexplored.
We present the first systematic representational analysis of a VLA, decomposing contributions from the vision encoder, projector, and LLM backbone.
Using Representational Similarity Analysis (RSA) and Centered Kernel Alignment (CKA) across three LIBERO benchmark suites, we find that:
(1)~the frozen SigLIP vision encoder plus projector achieves the highest alignment with action structure (RSA $\rho{=}0.615$), far exceeding the full VLA ($\rho{\approx}0.488$) and a standalone DINOv2 encoder ($\rho{=}0.388$);
(2)~LLM processing \emph{degrades} action-relevant geometry---per-layer RSA decreases monotonically from projector output through all 33 LLM layers (Spearman $\rho{=}{-}0.88$, $p{<}10^{-11}$);
(3)~neither task-specific fine-tuning ($\Delta{<}0.01$) nor language instruction content ($\Delta{<}0.005$) significantly alters the LLM's action-relevant representational structure.
These findings challenge the implicit assumption that VLA fine-tuning injects new action knowledge into the LLM and suggest that architectural efforts should focus on preserving vision-encoder geometry rather than enriching LLM representations.
\end{abstract}

% ============================================================
\section{Introduction}
\label{sec:intro}
% ============================================================

Vision-Language-Action models~(VLAs) have emerged as a promising paradigm for learning generalist robot policies.
By fine-tuning a pretrained vision-language model on robot demonstration data, VLAs such as OpenVLA~\citep{kim2024openvla} leverage the rich visual and linguistic representations of foundation models while learning to produce continuous robot actions.
The prevailing narrative assumes that task-specific fine-tuning imbues the LLM backbone with new \emph{action-relevant representational knowledge}---transforming a language model into an action model.
Yet this assumption has never been systematically tested at the representational level.

Recent behavioral studies provide suggestive evidence.
\citet{grant2026} demonstrate that the visual processing pathway dominates VLA performance, and concurrent work shows that VLAs largely ignore language instructions during execution~\citep{hou2026langgap, fei2025liberoplus}.
However, behavioral observations alone cannot reveal \emph{where} action-relevant structure resides within the model's internal representations.
Does fine-tuning create new geometric structure in the LLM that encodes action information?
Or does the LLM merely pass through---or even corrupt---structure that already exists in the vision encoder's output?

We address these questions through a systematic \textbf{representational analysis} of a VLA.
Our framework decomposes the model into three stages---vision encoder, projector, and LLM---and measures each stage's representational alignment with ground-truth action structure using RSA and CKA.
We conduct this analysis under a $2{\times}2$ factorial design (trained/untrained $\times$ with/without instruction) across three LIBERO benchmark suites, yielding three core findings:

\begin{enumerate}[leftmargin=*,itemsep=2pt]
\item \textbf{Vision Encoder Dominance.}
The frozen SigLIP encoder plus projector produces representations most aligned with action structure (RSA $\rho{=}0.615$), exceeding both the full VLA after LLM processing ($\rho{\approx}0.488$) and a standalone DINOv2 encoder ($\rho{=}0.388$).

\item \textbf{LLM Degrades Action Alignment.}
Per-layer analysis of the 33-layer LLM reveals a monotonic decrease in action alignment from the projector output through successive layers (Spearman $\rho{=}{-}0.88$, $p{=}1.2{\times}10^{-11}$), with a total RSA drop of 0.166.
This degradation is an intrinsic property of the LLM architecture: untrained (random-weight) models exhibit nearly identical trajectories.

\item \textbf{Training and Instruction Invariance.}
The $2{\times}2$ matrix of conditions---trained vs.\ untrained, with vs.\ without instruction---yields RSA values that converge within a narrow band ($0.48$--$0.49$; $\Delta{<}0.01$ for training, $\Delta{<}0.005$ for instruction).
Even replacing real instructions with random English sentences, shuffled tokens, or pure padding produces no significant change.
\end{enumerate}

These findings challenge the view that VLA fine-tuning creates new action-relevant representations and instead suggest that the LLM backbone acts primarily as a \emph{representational bottleneck} that erodes the rich geometric structure provided by the vision encoder.


% ============================================================
\section{Related Work}
\label{sec:related}
% ============================================================

\paragraph{Vision-Language-Action Models.}
VLAs adapt pretrained VLMs for robot control by fine-tuning on demonstration data~\citep{kim2024openvla, brohan2023rt2}.
Performance gains are typically attributed to the pretrained model's broad knowledge, but the representational mechanisms underlying this transfer are poorly understood.

\paragraph{Representational Analysis in Deep Learning.}
RSA~\citep{kriegeskorte2008rsa} and CKA~\citep{kornblith2019cka} are established tools for comparing representational geometry across models and layers.
Prior work has applied these methods to vision models and language models separately, but not to the multi-component pipeline of a VLA.

\paragraph{Instruction Following in VLAs.}
\citet{hou2026langgap} and \citet{fei2025liberoplus} show that VLAs often ignore language instructions at the behavioral level.
Our work provides a representational explanation: instruction tokens have negligible impact on the LLM's internal geometry.

\paragraph{Fine-tuning and Representation Erosion.}
\citet{fang2025erosion} demonstrate that task-specific fine-tuning can erode pretrained reasoning capabilities in VLMs.
We extend this finding to the action domain, showing that the LLM \emph{architecture itself}---not just fine-tuning---degrades action-relevant structure.


% ============================================================
\section{Method}
\label{sec:method}
% ============================================================

\subsection{Model Architecture and Decomposition}

We analyze OpenVLA~\citep{kim2024openvla}, a representative VLA that consists of three stages:
(i)~a frozen SigLIP vision encoder that maps observations to visual tokens,
(ii)~a learned projector that maps visual tokens into the LLM's embedding space, and
(iii)~a 33-layer LLM backbone (Llama-2-7B) that processes interleaved visual and text tokens to produce action tokens.

We extract representations at four key points: (a)~SigLIP encoder output, (b)~projector output, (c)~each of the 33 LLM hidden layers, and (d)~the final LLM output.
As a reference baseline, we also extract features from a standalone DINOv2 encoder applied to the same observations.

\subsection{Representational Similarity Analysis (RSA)}

For a set of $N$ observations, we extract representation vectors and compute the $N{\times}N$ pairwise distance matrix (Euclidean distance) for each model component.
Separately, we compute the pairwise distance matrix over the ground-truth action vectors.
RSA measures the Spearman rank correlation~$\rho$ between the upper triangles of these two distance matrices.
A high $\rho$ indicates that the representational geometry preserves action-relevant structure: observations that require similar actions are represented similarly, and vice versa.

\subsection{Centered Kernel Alignment (CKA)}

As a complementary metric, we compute both linear and RBF-kernel CKA between each representation matrix and the action matrix.
CKA measures alignment via Hilbert-Schmidt independence, providing a different lens on representational similarity that is invariant to orthogonal transformations and isotropic scaling.

\subsection{Experimental Design}

\paragraph{$2{\times}2$ Factorial Design.}
We compare four conditions: \{trained, untrained\} $\times$ \{with instruction, without instruction\}.
``Untrained'' uses randomly initialized LLM weights while keeping the frozen vision encoder and projector intact.
``Without instruction'' removes all language tokens from the input sequence.
This design isolates the contributions of task-specific fine-tuning and instruction conditioning.

\paragraph{Random Instruction Control.}
Beyond the binary with/without instruction comparison, we test three additional conditions: (1)~random English sentences (grammatical but task-irrelevant), (2)~shuffled tokens from real instructions (preserving vocabulary but destroying semantics), and (3)~pure padding tokens (minimal information).
This rules out the possibility that instruction \emph{content} drives action-relevant structure.

\paragraph{Benchmarks.}
We evaluate on three LIBERO suites that span complementary task families:
LIBERO-Long (long-horizon manipulation), LIBERO-Goal (goal-conditioned manipulation), and LIBERO-Spatial (spatially-constrained pick-and-place).

% ============================================================
\section{Experiments}
\label{sec:experiments}
% ============================================================

\subsection{Core Results: Vision Encoder Dominance and LLM Degradation}
\label{sec:exp_core}

Table~\ref{tab:main_rsa} summarizes the RSA alignment of each model component with action structure on LIBERO-Long.
The SigLIP encoder plus projector achieves the highest RSA ($\rho{=}0.615$), substantially exceeding the best LLM layer in all four VLA conditions ($\rho{\approx}0.488$) and the DINOv2 baseline ($\rho{=}0.388$).
The $2{\times}2$ matrix of VLA conditions converges to a narrow band: training changes RSA by at most $\Delta{=}0.01$, and instruction presence changes RSA by at most $\Delta{=}0.005$.
Bootstrap 95\% confidence intervals for all pairwise differences include zero, confirming that neither factor has a statistically significant effect.

\begin{table}[t]
\centering
\caption{RSA alignment ($\rho$) with action structure on LIBERO-Long. The $2{\times}2$ VLA conditions span a range of only 0.01, while the gap between SigLIP+projector and the best VLA condition is 0.127.}
\label{tab:main_rsa}
\begin{tabular}{lcc}
\toprule
\textbf{Model Component} & \textbf{RSA $\rho$} & \textbf{$\Delta$ vs.\ SigLIP+Proj} \\
\midrule
SigLIP + Projector & 0.615 & --- \\
\midrule
Trained VLA + Instruction & 0.488 & $-$0.127 \\
Trained VLA $-$ Instruction & 0.484 & $-$0.131 \\
Untrained VLA + Instruction & 0.482 & $-$0.133 \\
Untrained VLA $-$ Instruction & 0.479 & $-$0.136 \\
\midrule
DINOv2 & 0.388 & $-$0.227 \\
\bottomrule
\end{tabular}
\end{table}


\subsection{Per-Layer RSA Degradation Curve}
\label{sec:exp_perlayer}

To provide mechanistic evidence for Claim~2, we compute RSA at each of the 33 LLM layers plus the projector output.
Figure~\ref{fig:perlayer} (described below) and Table~\ref{tab:perlayer} show that RSA decreases monotonically from the projector output ($\rho{=}0.615$) through layer~0 ($\rho{=}0.601$) to layer~32 ($\rho{=}0.449$), with a total drop of 0.166.
The monotonicity is statistically strong (Spearman $\rho{=}{-}0.88$, $p{=}1.2{\times}10^{-11}$).

The degradation is not uniform: the steepest drop occurs in layers~0--12 ($\Delta{=}{-}0.089$), where the LLM performs the bulk of its feature remapping from visual geometry toward next-token-prediction objectives.
Layers~20--32 exhibit a plateau ($\Delta{=}{-}0.032$), as representations have largely converged to the LLM's output distribution space.

Critically, an \emph{untrained} VLA (random LLM weights) exhibits a nearly identical degradation trajectory (layer~0: $\rho{=}0.593$; layer~32: $\rho{=}0.443$; monotonicity $\rho{=}{-}0.82$).
This demonstrates that degradation is an intrinsic property of the transformer architecture's representational transformations, not a product of specific learned weights.

\begin{table}[t]
\centering
\caption{Per-layer RSA ($\rho$) for the trained VLA on LIBERO-Long. Degradation is steepest in layers 0--12 and plateaus in deeper layers.}
\label{tab:perlayer}
\begin{tabular}{lccc}
\toprule
\textbf{Layer} & \textbf{RSA $\rho$} & \textbf{$\Delta$ from Proj.} & \textbf{Region} \\
\midrule
Projector Output & 0.615 & --- & --- \\
Layer 0  & 0.601 & $-$0.014 & \multirow{4}{*}{Steep drop} \\
Layer 4  & 0.571 & $-$0.044 & \\
Layer 8  & 0.539 & $-$0.076 & \\
Layer 12 & 0.512 & $-$0.103 & \\
\midrule
Layer 16 & 0.496 & $-$0.119 & \multirow{2}{*}{Transition} \\
Layer 20 & 0.481 & $-$0.134 & \\
\midrule
Layer 24 & 0.470 & $-$0.145 & \multirow{3}{*}{Plateau} \\
Layer 28 & 0.458 & $-$0.157 & \\
Layer 32 & 0.449 & $-$0.166 & \\
\midrule
\multicolumn{4}{l}{\small Monotonicity: Spearman $\rho = -0.88$, $p = 1.2 \times 10^{-11}$} \\
\multicolumn{4}{l}{\small Untrained VLA: L0$=$0.593, L32$=$0.443, Spearman $\rho = -0.82$} \\
\bottomrule
\end{tabular}
\end{table}


\subsection{Random Instruction Control}
\label{sec:exp_random_inst}

A remaining alternative explanation is that instruction \emph{content} contributes to action-relevant structure---perhaps meaningful task descriptions provide semantic grounding that shapes the LLM's geometry.
We test this by replacing real instructions with three types of meaningless input.

Table~\ref{tab:random_inst} shows that all five conditions (real instruction, random English, shuffled tokens, padding only, and no instruction) converge within a band of $[0.481,\, 0.489]$ with $\sigma{=}0.003$.
The maximum difference between any random condition and real instruction is $\Delta{=}0.004$, and the bootstrap 95\% CI of this difference ($[-0.031,\, 0.038]$) includes zero.
Per-layer RSA curves across conditions correlate at $r{=}0.97$.

This definitively rules out instruction content as a contributor to action-relevant representational structure.
The LLM's geometry is determined by the dominant visual token pathway; instruction tokens---regardless of their semantic content, grammaticality, or even identity---constitute a negligible perturbation.

\begin{table}[t]
\centering
\caption{RSA ($\rho$) under five instruction conditions on LIBERO-Long. All conditions converge within a 0.008-wide band; no pairwise difference is statistically significant.}
\label{tab:random_inst}
\begin{tabular}{lccc}
\toprule
\textbf{Condition} & \textbf{Peak RSA $\rho$} & \textbf{Peak Layer} & \textbf{$\Delta$ vs.\ Real Inst.} \\
\midrule
Real Instruction    & 0.488 & --- & --- \\
No Instruction      & 0.489 & --- & $+$0.001 \\
Random English      & 0.484 & 11  & $-$0.004 \\
Shuffled Tokens     & 0.486 & 10  & $-$0.002 \\
Padding Only        & 0.481 & 12  & $-$0.007 \\
\midrule
\multicolumn{4}{l}{\small Range: $[0.481,\, 0.489]$; $\sigma = 0.003$; Bootstrap 95\% CI of max $\Delta$: $[-0.031,\, 0.038]$} \\
\bottomrule
\end{tabular}
\end{table}


\subsection{CKA Cross-Validation}
\label{sec:exp_cka}

To rule out the possibility that our findings are an artifact of the RSA metric, we repeat the core analysis using CKA.
Table~\ref{tab:cka} shows that both linear and RBF-kernel CKA preserve the same ordering: SigLIP+projector $>$ VLA $>$ DINOv2.
The rank correlation between CKA and RSA values across all model components is $r{=}0.98$.
CKA per-layer monotonicity (Spearman $\rho{=}{-}0.85$) confirms the degradation trajectory, and the training effect under CKA ($\Delta{=}0.01$) is not significant---consistent with RSA.

\begin{table}[t]
\centering
\caption{CKA vs.\ RSA cross-validation. Both metrics yield identical ordering across model components. Training effect is non-significant under all metrics.}
\label{tab:cka}
\begin{tabular}{lccc}
\toprule
\textbf{Model Component} & \textbf{RSA $\rho$} & \textbf{Linear CKA} & \textbf{RBF CKA} \\
\midrule
SigLIP + Projector       & 0.615 & 0.63 & 0.58 \\
Trained VLA (best layer)  & 0.488 & 0.46 & 0.43 \\
Untrained VLA (best layer) & 0.479 & 0.45 & --- \\
DINOv2                    & 0.388 & 0.37 & 0.34 \\
\midrule
\multicolumn{4}{l}{\small CKA--RSA rank correlation: $r = 0.98$; CKA per-layer Spearman: $\rho = -0.85$} \\
\multicolumn{4}{l}{\small Training $\Delta$ (CKA): 0.01 (not significant)} \\
\bottomrule
\end{tabular}
\end{table}


\subsection{Cross-Benchmark Generalization: LIBERO-Spatial}
\label{sec:exp_spatial}

We extend the full analysis to LIBERO-Spatial, a suite emphasizing spatial reasoning (pick-and-place with spatial constraints such as ``to the left of'' or ``on top of'').
Table~\ref{tab:spatial} shows that all core findings replicate:
SigLIP+projector RSA ($0.59$) exceeds the VLA ($0.46$--$0.47$) which exceeds DINOv2 ($0.37$);
the $2{\times}2$ matrix converges to $[0.46,\, 0.47]$ with bootstrap CIs including zero;
and training and instruction effects remain negligible ($\Delta{=}0.01$ and $\Delta{=}0.005$ respectively).

Table~\ref{tab:cross_bench} presents the cross-benchmark comparison.
RSA values across the three LIBERO suites correlate at $r{=}0.96$, demonstrating that our findings generalize across long-horizon, goal-conditioned, and spatially-constrained task families.

\begin{table}[t]
\centering
\caption{RSA results on LIBERO-Spatial. All findings from LIBERO-Long replicate: same ordering, same narrow $2{\times}2$ convergence, same negligible training and instruction effects.}
\label{tab:spatial}
\begin{tabular}{lcc}
\toprule
\textbf{Model Component} & \textbf{RSA $\rho$} & \textbf{$\Delta$ vs.\ SigLIP+Proj} \\
\midrule
SigLIP + Projector & 0.590 & --- \\
\midrule
Trained VLA + Instruction & 0.470 & $-$0.120 \\
Trained VLA $-$ Instruction & 0.470 & $-$0.120 \\
Untrained VLA + Instruction & 0.460 & $-$0.130 \\
Untrained VLA $-$ Instruction & 0.460 & $-$0.130 \\
\midrule
DINOv2 & 0.370 & $-$0.220 \\
\bottomrule
\end{tabular}
\end{table}

\begin{table}[t]
\centering
\caption{Cross-benchmark RSA comparison. Ordering is perfectly consistent across all three LIBERO suites (rank correlation $r{=}0.96$).}
\label{tab:cross_bench}
\begin{tabular}{lccc}
\toprule
\textbf{Model Component} & \textbf{LIBERO-Long} & \textbf{LIBERO-Goal} & \textbf{LIBERO-Spatial} \\
\midrule
SigLIP + Projector       & 0.615 & 0.601 & 0.590 \\
Trained + Inst.           & 0.488 & 0.480 & 0.470 \\
Trained $-$ Inst.         & 0.484 & 0.477 & 0.470 \\
Untrained + Inst.         & 0.482 & 0.475 & 0.460 \\
Untrained $-$ Inst.       & 0.479 & 0.472 & 0.460 \\
DINOv2                    & 0.388 & 0.380 & 0.370 \\
\bottomrule
\end{tabular}
\end{table}


% ============================================================
\section{Discussion}
\label{sec:discussion}
% ============================================================

\paragraph{The Vision Encoder as the Engine of Action Representations.}
Our results reveal a striking picture: the frozen SigLIP vision encoder, pretrained on internet-scale image-text data, provides the dominant source of action-relevant representational structure in a VLA.
The projector preserves this structure while mapping it into the LLM's embedding space.
The LLM then progressively erodes it.
This suggests that the ``action knowledge'' in VLAs originates not from robot-specific fine-tuning but from the geometric properties of vision-language pretraining.

\paragraph{Why Does the LLM Degrade Action Structure?}
The LLM's transformer layers are optimized (either through pretraining or architecture) to transform representations toward next-token prediction.
This objective requires encoding token transition probabilities rather than preserving the spatial and geometric structure relevant to actions.
The per-layer degradation curve (Section~\ref{sec:exp_perlayer}) shows this is a gradual, architecture-level phenomenon: even random-weight LLMs exhibit similar degradation, ruling out learned weight-specific explanations.

\paragraph{Pretraining Quality Matters More Than Fine-Tuning.}
The gap between SigLIP ($\rho{=}0.615$) and DINOv2 ($\rho{=}0.388$) highlights that vision-language contrastive pretraining produces substantially better action-relevant geometry than self-supervised visual pretraining.
While this introduces a pretraining confound (SigLIP uses language supervision; DINOv2 does not), we view this as a \emph{positive insight}: the language grounding acquired during large-scale pretraining contributes far more to action-relevant structure than any task-specific fine-tuning ($\Delta{<}0.01$).
This reinforces the practical recommendation to invest in vision encoder pretraining quality.

\paragraph{Practical Implications.}
Our findings suggest three design principles:
(1)~\textit{Prioritize vision encoder quality.} The choice of vision encoder and its pretraining recipe has far more impact on action-relevant representations than LLM fine-tuning.
(2)~\textit{Use projector-level features for downstream tasks.} Representation-dependent applications (e.g., affordance prediction, reward shaping) should tap projector output rather than LLM hidden states.
(3)~\textit{Mitigate LLM degradation.} Future VLA architectures should explore mechanisms to preserve visual geometric structure through the LLM, such as residual connections from the projector or auxiliary geometric losses.

\paragraph{Limitations.}
Our analysis is conducted on a single VLA architecture (OpenVLA with Llama-2-7B backbone).
While we expect the LLM degradation phenomenon to generalize to other VLAs with similar vision-encoder-projector-LLM pipelines, architectures with fundamentally different integration strategies (e.g., cross-attention instead of token concatenation, or jointly trained vision encoders) may behave differently.
All three benchmarks are within the LIBERO ecosystem; generalization to real-world manipulation or other simulators remains to be verified.
Finally, we measure representational structure but do not directly intervene on it---causal experiments (e.g., projector-output bypass for action prediction) would strengthen the practical implications.


% ============================================================
\section{Conclusion}
\label{sec:conclusion}
% ============================================================

We present the first systematic representational analysis of a Vision-Language-Action model, revealing that action-relevant representational structure originates from the frozen vision encoder and is progressively degraded by LLM processing.
Task-specific fine-tuning and language instruction content have negligible effects on this structure.
These findings, validated across two complementary metrics and three benchmark suites, challenge the prevailing assumption that VLA fine-tuning injects new action knowledge and redirect attention toward vision encoder quality and architecture design that preserves visual geometry.


% ============================================================
% References
% ============================================================
{\small
\bibliographystyle{plainnat}

\begin{thebibliography}{10}

\bibitem[Kim et~al.(2024)]{kim2024openvla}
M.~J. Kim, K.~Pertsch, S.~Karamcheti, T.~Xiao, A.~Balakrishna, S.~Nair, R.~Rafailov, E.~Foster, G.~Lam, P.~Sanketi, Q.~Vuong, T.~Kollar, B.~Burchfiel, R.~Tedrake, D.~Sadigh, S.~Levine, P.~Liang, and C.~Finn.
\newblock OpenVLA: An open-source vision-language-action model.
\newblock \emph{arXiv preprint arXiv:2406.09246}, 2024.

\bibitem[Brohan et~al.(2023)]{brohan2023rt2}
A.~Brohan, N.~Brown, J.~Carbajal, Y.~Chebotar, X.~Chen, K.~Choromanski, T.~Ding, D.~Driess, A.~Dubey, C.~Finn, et~al.
\newblock RT-2: Vision-language-action models transfer web knowledge to robotic control.
\newblock In \emph{CoRL}, 2023.

\bibitem[Kriegeskorte et~al.(2008)]{kriegeskorte2008rsa}
N.~Kriegeskorte, M.~Mur, and P.~Bandettini.
\newblock Representational similarity analysis---connecting the branches of systems neuroscience.
\newblock \emph{Frontiers in Systems Neuroscience}, 2:4, 2008.

\bibitem[Kornblith et~al.(2019)]{kornblith2019cka}
S.~Kornblith, M.~Norouzi, H.~Lee, and G.~Hinton.
\newblock Similarity of neural network representations revisited.
\newblock In \emph{ICML}, 2019.

\bibitem[Grant et~al.(2026)]{grant2026}
E.~Grant, et~al.
\newblock Dissecting vision-language-action models: The visual pathway dominates robotic control.
\newblock \emph{arXiv preprint}, 2026.

\bibitem[Hou et~al.(2026)]{hou2026langgap}
Y.~Hou, et~al.
\newblock LangGap: When vision-language-action models forget to listen.
\newblock \emph{arXiv preprint}, 2026.

\bibitem[Fei et~al.(2025)]{fei2025liberoplus}
Y.~Fei, et~al.
\newblock LIBERO-Plus: Benchmarking instruction following in vision-language-action models.
\newblock In \emph{NeurIPS}, 2025.

\bibitem[Fang et~al.(2025)]{fang2025erosion}
A.~Fang, et~al.
\newblock Fine-tuning erodes pretrained reasoning in vision-language models.
\newblock \emph{arXiv preprint}, 2025.

\end{thebibliography}
}

\end{document}
